\documentclass[11pt]{article}
\usepackage[margin=1in]{geometry}
\usepackage{amsmath,amssymb}
\title{Tax-Optimal Withdrawals Under FIFO Lot Accounting:\\ An Exact Optimisation Approach for Indian Listed Equity}
\author{Yashasvi Adlak}
\date{}
\begin{document}
\maketitle

\begin{abstract}
% TODO write last. Problem, method (MILP), main result (tax saved vs baselines), LLM comparison.
\end{abstract}

\section{Introduction}
% TODO: why withdrawals trigger avoidable tax; why retail tools (broker P\&L) only report, not plan;
% why FIFO changes the problem; contributions (3 bullets).

\section{Legal Model}
% TODO cite primary sources. Holding period (more than 12 months), rates, exemption,
% set-off rules, FIFO per demat account (s.45(2A), 1961 Act; new Act section TBD).

\section{Problem Formulation}
A \emph{position} $p$ is one scrip in one demat account with price $P_p$ and FIFO-ordered lots $j=1,\dots,n_p$
with quantity $q_{pj}$ and unit cost $c_{pj}$. The decision is the integer quantity $x_p \in \{0,\dots,Q_p\}$ sold,
where $Q_p=\sum_j q_{pj}$. FIFO fixes the consumed quantity $y_{pj}(x_p)$ of each lot, so lot choice is not a variable.

Realised gains: $G^{ST}=\sum_{p,j\in S}(P_p-c_{pj})y_{pj}$ and $G^{LT}=\sum_{p,j\in L}(P_p-c_{pj})y_{pj}$,
where a lot is long-term iff held more than 12 months at the sale date. Adding gains already realised this year,
a net short-term loss offsets long-term gains but a net long-term loss does not offset short-term gains:
\[
T(G^{ST},G^{LT}) = (1+\kappa)\Big(\tau_s \max(\tilde G^{ST},0) + \tau_l \max(\max(\tilde G^{LT},0)-E,\,0)\Big),
\]
with $\tilde G^{LT}=G^{LT}+\min(G^{ST},0)$ when $G^{ST}<0$ and $G^{LT}>0$ (and $\tilde G^{LT}=G^{LT}$ otherwise),
$\tilde G^{ST}=\max(G^{ST},0)$, rates $\tau_s=0.20$, $\tau_l=0.125$, cess $\kappa=0.04$, exemption $E=\text{Rs }125{,}000$.

\textbf{Problem.} $\min_x\; T(G^{ST}(x),G^{LT}(x))$ subject to $\sum_p P_p x_p \ge W$, $x_p\in\{0,\dots,Q_p\}$.

\section{Method}
We solve the problem exactly with CP-SAT over integer paise. For each position, lots are ordered by purchase date and the
variable $y_{pj}\in[0,q_{pj}]$ is the quantity taken from lot $j$. A binary $z_{pj}$ marks lot $j$ as fully consumed, with
$y_{pj}\ge q_{pj}z_{pj}$ and $y_{p,j+1}\le q_{p,j+1}z_{pj}$, which enforces FIFO order. Gains are linear in $y$.
The tax function is encoded with min/max equalities: $u=\min(\max(-\tilde G^{ST},0),\max(\tilde G^{LT},0))$ for the
short-term-loss set-off, further brought-forward set-offs, $\min$ for the exemption, and a pair of linear inequalities
that reproduce round-half-up to the paisa. Because a single weighted objective overflows 64-bit integers on large
portfolios, we solve lexicographically in two stages: minimise tax, then fix it and minimise proceeds.
Correctness is validated against exhaustive search (80 random instances with random year-to-date gains and
brought-forward losses; identical tax and proceeds on all).

\section{Baselines}
\textbf{Pro-rata:} sell the same fraction of every position (floor of the target, then top up the largest deficits).
\textbf{Largest-first:} empty the biggest positions until the target is met.
\textbf{Greedy tax-aware:} repeatedly sell the chunk with the lowest marginal tax per rupee, evaluated through the full FIFO and tax model; strong but not guaranteed optimal.
\textbf{Exhaustive optimum:} enumerate every feasible quantity vector (small instances only); ground truth for the exact solver.
% TODO LLM-generated plans (later).

\section{Experiments and Results}
\textbf{Instances.} Synthetic, seeded, deliberately hard: gains-heavy portfolios (cost 0.3--1.15 of price), 40\% of lots acquired
10--14 months before sale (near the long-term boundary), plus random year-to-date gains and brought-forward losses.
Withdrawals are 20--50\% of portfolio value. Sizes: 15, 48 and 100 lots (100, 40 and 15 instances).
\textbf{Validation.} On 80 small random instances the solver matches exhaustive search exactly (tax and proceeds); on the
hard instances no baseline ever beat it and a single-worker re-solve agreed on 25 of 25 checked instances.
\textbf{Results.} Table~\ref{tab:main} reports mean excess tax over the solver's optimum.
\begin{table}[h]\centering\begin{tabular}{rrrrrrr}
\hline
Lots & Inst. & Solver (s) & Greedy subopt.\% & Greedy excess (Rs) & Pro-rata excess (Rs) & Largest-first excess (Rs) \\
\hline
15 & 100 & 0.02 & 87.0 & 14,235 & 58,552 & 73,573 \\
48 & 40 & 0.07 & 92.5 & 57,490 & 238,622 & 315,717 \\
100 & 15 & 0.15 & 93.3 & 95,905 & 523,565 & 636,212 \\
\hline
\end{tabular}
\caption{Baselines versus the exact solver on hard synthetic instances.}\label{tab:main}\end{table}
The solver proved optimality in every instance (mean 0.02--0.15\,s). The greedy heuristic was suboptimal in 87--93\% of instances,
pro-rata and largest-first in 96--100\%. \emph{Caveat:} greedy here is our own marginal-tax heuristic with coarse chunking;
a stronger rule-based practitioner baseline should be added before claiming general superiority.
% TODO LLM comparison: zero-shot, CoT, code interpreter, solver-as-tool.

\section{Limitations}
Single financial year; listed equity only; no surcharge, brokerage or STT; prices treated as known; synthetic instances only
(no real client portfolios yet); several rule parameters pending primary-source verification; ingestion ignores corporate actions.

\section{Conclusion}
\end{document}
